\newcommand{\DateField}[3]{
\begin{ocg}{ocg#1}{ocg#1}{false}
\begin{tikzpicture}[overlay, remember picture]
    \draw[line width=1pt, red]
    (-0.00, -0.25) rectangle (3.20, 0.80); % Größe des Rahmens
\end{tikzpicture}
\end{ocg}
\TextField[
    onblur={dateChecker(#2, #3);}, % range in month (negativ are in past, positiv are in futute)
    validate={dateChecker(#2, #3);}, % range in month (negativ are in past, positiv are in futute)
    name=#1,
    format={AFDate_FormatEx("dd.mm.yyyy");},
    keystroke={AFDate_KeystrokeEx("dd.mm.yyyy");},
    bordercolor={0 0 0},
    borderwidth=1pt,
    borderstyle=S,
    height=0.7cm,
    width=3cm,
]{}
}

\newcommand{\EmailField}[1]{
\begin{ocg}{ocg#1}{ocg#1}{false}
\begin{tikzpicture}[overlay, remember picture]
    \draw[line width=1pt, red]
    (-0.00, -0.25) rectangle (10.20, 0.80); % Größe des Rahmens
\end{tikzpicture}
\end{ocg}
\TextField[
    name=#1,
    keystroke={emailKeystroke();},
    onblur={emailChecker();},
    validate={emailChecker();},
    bordercolor={0 0 0},
    width=10cm,
    height=0.7cm
]{}
}

\newcommand{\PhoneField}[1]{
\begin{ocg}{ocg#1}{ocg#1}{false}
\begin{tikzpicture}[overlay, remember picture]
    \draw[line width=1pt, red]
    (-0.00, -0.25) rectangle (10.20, 0.80); % Größe des Rahmens
\end{tikzpicture}
\end{ocg}
\TextField[
    name=#1,
    keystroke={phoneKeystroke();},
    onblur={phoneChecker();},
    validate={phoneChecker();},
    bordercolor={0 0 0},
    width=10cm,
    height=0.7cm
]{}
}

\newcommand{\IBANField}[1]{
\begin{ocg}{ocg#1}{ocg#1}{false}
\begin{tikzpicture}[overlay, remember picture]
    \draw[line width=1pt, red]
    (-0.00, -0.25) rectangle (8.20, 0.80); % Größe des Rahmens
\end{tikzpicture}
\end{ocg}
\TextField[
    name=#1,
    keystroke={ibanKeystroke();},
    onblur={ibanChecker();},
    validate={ibanChecker();},
    format={var f = this.getField(event.targetName); f.textFont="Courier";},
    bordercolor={0 0 0},
    width=8cm,
    height=0.7cm
]{}
}